\section{Score Matching}

上一章中，我们通过 Flow Matching 学习了一个向量场 $u_t^\theta$，使得对应 ODE 能够把简单分布 $p_{init}$ 推到数据分布 $\data$。这一章介绍另一条经典路线：不直接学习速度场，而是学习每个时刻的 \emph{score function}，再利用 score 来构造生成过程。

\subsection{条件 score 与边缘 score}
设 $p(x)$ 是 $\R^d$ 上的一个概率分布，则其 score function 定义为 $s_p(\vx)=\nabla_\vx \log p(\vx)$。


score 描述的是密度在当前位置增长最快的方向。由于
\(
\nabla_\vx \log p(\vx)=\frac{\nabla_\vx p(\vx)}{p(\vx)}
\)，
它可以理解为“相对密度变化率”。

在 Flow Matching 中，我们先指定条件概率路径 $p_t(\cdot|\vz)$，再通过边缘化得到边缘概率路径
\begin{equation*}
p_t(\vx)=\int p_t(\vx|\vz)\data(\vz)\,\dd \vz.
\end{equation*}
对于 score，边缘score可以通过条件score function 表示为
\begin{equation}
  \nabla_\vx \log p_t(\vx)
  =
  \int
  \nabla_\vx \log p_t(\vx|\vz)
  \frac{p_t(\vx|\vz)\data(\vz)}{p_t(\vx)}
  \,\dd \vz.
  \end{equation}

\begin{proof}
由边缘密度的定义，
\begin{equation*}
p_t(\vx)=\int p_t(\vx|\vz)\data(\vz)\,\dd \vz.
\end{equation*}
对 $\vx$ 求梯度得
\begin{equation*}
\nabla_\vx p_t(\vx)
=
\int \nabla_\vx p_t(\vx|\vz)\data(\vz)\,\dd \vz.
\end{equation*}
两边同时除以 $p_t(\vx)$，并利用
\(
\nabla_\vx p_t(\vx|\vz)
=
p_t(\vx|\vz)\nabla_\vx \log p_t(\vx|\vz)
\)
即可得到
\begin{equation*}
\nabla_\vx \log p_t(\vx)
=
\int
\nabla_\vx \log p_t(\vx|\vz)
\frac{p_t(\vx|\vz)\data(\vz)}{p_t(\vx)}
\,\dd \vz.
\end{equation*}
\end{proof}


\begin{example}[Gaussian 条件概率路径的条件 score]
考虑 Gaussian 条件概率路径
\begin{equation}
p_t(\cdot|\vz)=\Normal(\alpha_t \vz,\beta_t^2 I_d).
\end{equation}
其对数密度为
\begin{equation}
\log p_t(\vx|\vz)
=
-\frac{d}{2}\log(2\pi \beta_t^2)
-\frac{1}{2\beta_t^2}\|\vx-\alpha_t \vz\|^2.
\end{equation}
因此条件 score function为
\begin{equation}
\nabla_\vx \log p_t(\vx|\vz)
=
-\frac{\vx-\alpha_t \vz}{\beta_t^2}
\end{equation}
\end{example}
注意到 Gaussian 条件概率路径的条件 score function 是一个线性函数，且与 $\vz$ 的关系完全由均值 $\alpha_t \vz$ 决定，同时 条件向量场 $u_t(x|z)$ 也具有该性质。因此二者之间存在转换关系

\begin{proposition}[Gaussian 概率路径下向量场与 score 的转换公式]
\label{prop:convertion}
  对于 Gaussian 条件概率路径
\begin{equation}
p_t(\cdot|\vz)=\Normal(\alpha_t \vz,\beta_t^2 I_d),
\end{equation}
条件（或边缘）向量场与条件（或边缘） score 满足
\begin{equation}
  \begin{aligned}
    u_t^{\text{target}}(x|z) &= a_t \nabla \log p_t(x|z) + b_t x, \quad a_t = \left( \beta_t^2 \frac{\dot{\alpha}_t}{\alpha_t} - \dot{\beta}_t \beta_t \right), \quad b_t = \frac{\dot{\alpha}_t}{\alpha_t} \\
    u_t^{\text{target}}(x) &= a_t \nabla \log p_t(x) + b_t x. 
  \end{aligned}
\end{equation}
\end{proposition}

\begin{proof}
从条件向量场和条件score，我们可以得到：
\begin{equation}
\begin{aligned}
u_t^{\text{target}}(x|z) &= \left( \dot{\alpha}_t - \frac{\dot{\beta}_t}{\beta_t} \alpha_t \right) z + \frac{\dot{\beta}_t}{\beta_t} x \\
&\stackrel{(i)}{=} \left( \beta_t^2 \frac{\dot{\alpha}_t}{\alpha_t} - \dot{\beta}_t \beta_t \right) \left( \frac{\alpha_t z - x}{\beta_t^2} \right) + \frac{\dot{\alpha}_t}{\alpha_t} x \\
&= \left( \beta_t^2 \frac{\dot{\alpha}_t}{\alpha_t} - \dot{\beta}_t \beta_t \right) \nabla \log p_t(x|z) + \frac{\dot{\alpha}_t}{\alpha_t} x
\end{aligned}
\end{equation}
在 $(i)$ 中我们做了一些数学代换。通过积分，我们可以得到边缘向量场和边缘score之间的关系
\begin{equation}
\begin{aligned}
u_t^{\text{target}}(x) &= \int u_t^{\text{target}}(x|z) \frac{p_t(x|z)p_{\text{data}}(z)}{p_t(x)} dz \\
&= \int [a_t \nabla \log p_t(x|z) + b_t x] \frac{p_t(x|z)p_{\text{data}}(z)}{p_t(x)} dz \\
&= a_t \int \left( \frac{\nabla p_t(x|z)}{p_t(x|z)} \right) \frac{p_t(x|z)p_{\text{data}}(z)}{p_t(x)} dz + b_t x \int \frac{p_t(x, z)}{p_t(x)} dz \\
&= \frac{a_t}{p_t(x)} \int \nabla p_t(x|z) p_{\text{data}}(z) dz + b_t x \int p_t(z|x) dz \\
&\stackrel{(i)}{=} \frac{a_t}{p_t(x)} \nabla \int p_t(x|z) p_{\text{data}}(z) dz + b_t x \cdot 1 \\
&= a_t \frac{\nabla p_t(x)}{p_t(x)} + b_t x \\
&= a_t \nabla \log p_t(x) + b_t x
\end{aligned}
\end{equation}
其中，$(i)$ 首先利用莱布尼茨法则将 $\nabla$ 提到积分号外面；其次利用了条件分布积分归一化 $\int p_t(z|x) dz = 1$。
\end{proof}

\begin{remark}[score 的重参数化]
在 Gaussian probability path 下，速度场参数化与 score 参数化本质上是等价的线性重参数化。换句话说，只要能够恢复后验均值 $\E[\vz|\vx]$，就可以在向量场、score 和 denoiser 三种表述之间互相转换。

定义条件denoiser和边缘denoiser
\begin{equation}
D_t(x|z) = z, \quad D_t(x) = \int z \frac{p_t(x|z) p_{\text{data}}(z)}{p_t(x)} dz \stackrel{(i)}{=} \frac{1}{\dot{\alpha}_t \beta_t - \alpha_t \dot{\beta}_t} (\beta_t u_t^{\text{target}}(x_t) - \dot{\beta}_t x_t).
\end{equation}
其中，$(i)$ 由~\ref{prop:convertion}推导得出。
学习 denoiser 的模型被称为 denoising diffusion model，与学习向量场和 score 等价。
\end{remark}

\subsection{如何利用 score 进行采样}

上一节说明了：对于 Gaussian 概率路径，边缘向量场
$u_t^{\mathrm{target}}(\vx)$ 和边缘 score
$\nabla_\vx \log p_t(\vx)$ 可以相互转换。现在我们讨论如何用 score
把一个确定性的 flow model 扩展成带随机性的 diffusion model。

先回忆上一章的结论：如果边缘向量场 $u_t^{\mathrm{target}}$ 满足
\begin{equation}
  X_0\sim p_{\mathrm{init}},
  \qquad
  \dd X_t=u_t^{\mathrm{target}}(X_t)\,\dd t,
\end{equation}
那么 $X_t\sim p_t$，特别地 $X_1\sim \data$。也就是说，这个 ODE
沿着我们指定的概率路径 $p_t$ 把噪声分布推到数据分布。接下来要问的是：
能不能在这个动力系统中加入随机噪声，同时仍然保证每个时刻的边缘分布还是
$p_t$？

答案是可以的。关键工具是 Fokker--Planck 方程，它是连续性方程在 SDE
情形下的推广。我们首先通过下式定义 \emph{拉普拉斯（Laplacian）算子} $\Delta$：
\begin{equation}
\Delta w_t(\vx)
=
\sum_{i=1}^d \frac{\partial^2}{\partial x_i^2} w_t(\vx)
=
\divergence(\nabla_\vx w_t)(\vx),
\end{equation}
其中 $w_t:\R^d\to\R$ 是标量场。

\begin{theorem}[Fokker--Planck 方程]
考虑 SDE
\begin{equation}
  X_0\sim p_{\mathrm{init}},
  \qquad
  \dd X_t=u_t(X_t)\,\dd t+\sigma_t\,\dd W_t,
\end{equation}
其中 $\sigma_t\ge 0$ 是只依赖时间的 diffusion coefficient。如果
$X_t$ 的密度为 $p_t$，则 $p_t$ 满足
\begin{equation}
  \partial_t p_t(\vx)
  =
  -\divergence\!\left(p_t u_t\right)(\vx)
  +\frac{\sigma_t^2}{2}\Delta p_t(\vx).
\end{equation}
反过来，若一条概率路径 $p_t$ 满足上式和初值条件，则它就是该 SDE
的边缘密度路径。
\end{theorem}

当 $\sigma_t=0$ 时，Laplacian 项消失，Fokker--Planck 方程就退化为
flow model 中的连续性方程
\begin{equation}
  \partial_t p_t(\vx)
  =
  -\divergence\!\left(p_t u_t\right)(\vx).
\end{equation}
因此，SDE 和 ODE 的区别正体现在额外的扩散项
$\frac{\sigma_t^2}{2}\Delta p_t$ 上。为了让加入噪声后的过程仍然沿着同一条
概率路径 $p_t$ 运动，我们需要在 drift 中加入一个由 score 给出的修正项。

\begin{theorem}[SDE extension trick]
\label{thm:SDE_extension}
设 $u_t^{\mathrm{target}}$ 是对应概率路径 $p_t$ 的边缘向量场，即
\begin{equation}
  X_0\sim p_{\mathrm{init}},
  \qquad
  \dd X_t=u_t^{\mathrm{target}}(X_t)\,\dd t
  \quad\Longrightarrow\quad
  X_t\sim p_t.
\end{equation}
则对任意 diffusion coefficient $\sigma_t\ge 0$，SDE
\begin{equation}
  X_0\sim p_{\mathrm{init}},
  \qquad
  \dd X_t=
  \left[
    u_t^{\mathrm{target}}(X_t)
    +\frac{\sigma_t^2}{2}\nabla_\vx \log p_t(X_t)
  \right]\dd t
  +\sigma_t\,\dd W_t
\end{equation}
仍然满足
\begin{equation}
  X_t\sim p_t,\qquad 0\le t\le 1.
\end{equation}
特别地，$X_1\sim \data$。
\end{theorem}

\begin{proof}
因为 $u_t^{\mathrm{target}}$ 对应的 ODE 沿着概率路径 $p_t$ 运动，
所以 $p_t$ 满足连续性方程
\begin{equation}
  \partial_t p_t(\vx)
  =
  -\divergence\!\left(p_t u_t^{\mathrm{target}}\right)(\vx).
\end{equation}
在右侧同时加上并减去同一个 Laplacian 项：
\begin{equation}
\begin{aligned}
  \partial_t p_t(\vx)
  &=
  -\divergence\!\left(p_t u_t^{\mathrm{target}}\right)(\vx)
  -\frac{\sigma_t^2}{2}\Delta p_t(\vx)
  +\frac{\sigma_t^2}{2}\Delta p_t(\vx)  \\
  &=
  -\divergence\!\left(p_t u_t^{\mathrm{target}}\right)(\vx)
  -\divergence\!\left(
      p_t \frac{\sigma_t^2}{2}\nabla_\vx \log p_t
    \right)(\vx)
  +\frac{\sigma_t^2}{2}\Delta p_t(\vx)  \\
  &=
  -\divergence\!\left(
      p_t\left[
        u_t^{\mathrm{target}}
        +\frac{\sigma_t^2}{2}\nabla_\vx \log p_t
      \right]
    \right)(\vx)
  +\frac{\sigma_t^2}{2}\Delta p_t(\vx).
\end{aligned}
\end{equation}
这里用到了
\(
  p_t\nabla_\vx \log p_t=\nabla_\vx p_t
\)
以及
\(
  \Delta p_t=\divergence(\nabla_\vx p_t)
\)。
最后一行正是上面 SDE 的 Fokker--Planck 方程。因此该 SDE 的边缘密度仍为
$p_t$。
\end{proof}

这个结论的含义很直接：一旦我们知道边缘向量场
$u_t^{\mathrm{target}}$ 和边缘 score $\nabla_\vx \log p_t$，就可以在不改变
目标概率路径的情况下，任意选择 $\sigma_t$ 来加入随机性。当
$\sigma_t=0$ 时，它退化回原来的 flow model；当 $\sigma_t>0$ 时，轨迹会变成
随机的，但每个时刻的样本分布仍然是同一个 $p_t$。

\begin{example}[Gaussian 概率路径下的 SDE 采样]
对于 Gaussian 概率路径，由命题 \ref{prop:convertion} 可知
\begin{equation}
  u_t^{\mathrm{target}}(\vx)
  =
  a_t\nabla_\vx\log p_t(\vx)+b_t\vx.
\end{equation}
因此 SDE extension trick 可以写成只依赖 score 的形式：
\begin{equation}
  \dd X_t
  =
  \left[
    \left(a_t+\frac{\sigma_t^2}{2}\right)
    \nabla_\vx\log p_t(X_t)
    +b_t X_t
  \right]\dd t
  +\sigma_t\,\dd W_t.
\end{equation}
如果我们用神经网络 $s_t^\theta(\vx)$ 近似
$\nabla_\vx\log p_t(\vx)$，训练后即可模拟
\begin{equation}
  \dd X_t
  =
  \left[
    \left(a_t+\frac{\sigma_t^2}{2}\right)
    s_t^\theta(X_t)
    +b_t X_t
  \right]\dd t
  +\sigma_t\,\dd W_t,
  \qquad
  X_0\sim p_{\mathrm{init}},
\end{equation}
并把终点 $X_1$ 作为生成样本。
\end{example}


需要强调的是，理论上任意 $\sigma_t\ge 0$ 都保持同一条边缘概率路径；
但在实际模型中，神经网络存在近似误差，数值求解 SDE 也存在离散化误差。
因此不同的 $\sigma_t$ 会影响采样质量和计算成本，通常需要通过实验选择。

\subsection{学习 score：Score Matching}

现在的问题变成：如何学习真实 score $s_t(\vx)$？

最直接的目标是最小化
\begin{equation}
\mathcal{L}_{\mathrm{SM}}(\theta)
=
\E_{t\sim \mathrm{Unif}[0,1],\,\vx\sim p_t}
\left[
\left\|
s_t^\theta(\vx)-s_t(\vx)
\right\|^2
\right],
\end{equation}
其中 $s_t^\theta$ 是神经网络给出的 score 估计。

但和 Flow Matching 一样，边缘 score $s_t(\vx)$ 通常无法直接计算。于是我们转而考虑条件 score，并定义 Denoising Score Matching loss：
\begin{equation}
\mathcal{L}_{\mathrm{DSM}}(\theta)
=
\E_{t\sim \mathrm{Unif}[0,1],\,\vz\sim \data,\,\vx\sim p_t(\cdot|\vz)}
\left[
\left\|
s_t^\theta(\vx)-s_t^{\mathrm{cond}}(\vx|\vz)
\right\|^2
\right].
\end{equation}

\begin{theorem}[Denoising Score Matching]
Score Matching loss 与 Denoising Score Matching loss 只相差一个与参数 $\theta$ 无关的常数：
\begin{equation}
\mathcal{L}_{\mathrm{SM}}(\theta)
=
\mathcal{L}_{\mathrm{DSM}}(\theta)+C.
\end{equation}
因此
\begin{equation}
\nabla_\theta \mathcal{L}_{\mathrm{SM}}(\theta)
=
\nabla_\theta \mathcal{L}_{\mathrm{DSM}}(\theta).
\end{equation}
\end{theorem}

\begin{proof}
证明与上一章的 Conditional Flow Matching 完全平行。由
\begin{equation}
\mathcal{L}_{\mathrm{SM}}(\theta)
=
\E_{t,\vx\sim p_t}
\left[
\left\|
s_t^\theta(\vx)-s_t(\vx)
\right\|^2
\right]
\end{equation}
出发，展开平方项得
\begin{equation}
\begin{aligned}
\mathcal{L}_{\mathrm{SM}}(\theta)
&=
\E_{t,\vx\sim p_t}\!\left[\|s_t^\theta(\vx)\|^2\right]
-2\E_{t,\vx\sim p_t}\!\left[s_t^\theta(\vx)^\top s_t(\vx)\right]
+C_1.
\end{aligned}
\end{equation}
再利用边缘化技巧
\begin{equation}
s_t(\vx)
=
\int
s_t^{\mathrm{cond}}(\vx|\vz)
\frac{p_t(\vx|\vz)\data(\vz)}{p_t(\vx)}
\,\dd \vz,
\end{equation}
可将中间项改写为
\begin{equation}
\E_{t,\vx\sim p_t}\!\left[s_t^\theta(\vx)^\top s_t(\vx)\right]
=
\E_{t,\vz,\vx\sim p_t(\cdot|\vz)}
\left[
s_t^\theta(\vx)^\top s_t^{\mathrm{cond}}(\vx|\vz)
\right].
\end{equation}
于是
\begin{equation}
\mathcal{L}_{\mathrm{SM}}(\theta)
=
\E_{t,\vz,\vx\sim p_t(\cdot|\vz)}
\left[
\|s_t^\theta(\vx)\|^2
-2s_t^\theta(\vx)^\top s_t^{\mathrm{cond}}(\vx|\vz)
\right]
+C_1.
\end{equation}
再在期望中加上并减去
\(
\|s_t^{\mathrm{cond}}(\vx|\vz)\|^2
\)，
就得到
\begin{equation}
\mathcal{L}_{\mathrm{SM}}(\theta)
=
\mathcal{L}_{\mathrm{DSM}}(\theta)+C,
\end{equation}
其中常数项 $C$ 与参数 $\theta$ 无关。
\end{proof}

\begin{example}[Gaussian path 下的 Denoising Score Matching]
对于 Gaussian 条件概率路径，
\begin{equation}
\vx=\alpha_t\vz+\beta_t\epsilon,
\qquad
\epsilon\sim \Normal(0,I_d),
\end{equation}
且条件 score 满足
\begin{equation}
s_t^{\mathrm{cond}}(\vx|\vz)
=
-\frac{\epsilon}{\beta_t}.
\end{equation}
因此 DSM loss 可写成
\begin{equation}
\mathcal{L}_{\mathrm{DSM}}(\theta)
=
\E_{t,\vz,\epsilon}
\left[
\left\|
s_t^\theta(\alpha_t\vz+\beta_t\epsilon)
+\frac{\epsilon}{\beta_t}
\right\|^2
\right].
\end{equation}
\end{example}

这正是现代扩散模型中最常见的训练目标之一。网络输入一个带噪样本和时间 $t$，输出对应时刻的 score；训练完成后，再用上一节的反向 SDE 从噪声逐步采样回数据。

\begin{summarybox}
Score-based generative modeling 的核心逻辑与 Flow Matching 非常相似：先指定一条从数据分布到简单分布的随机扰动路径，再学习与这条路径对应的目标函数。不同之处在于，这里学习的对象是 score
\(
s_t(\vx)=\nabla_\vx \log p_t(\vx)
\)，
而不是速度场。

对于条件概率路径 $p_t(\cdot|\vz)$，我们往往可以显式写出条件 score $s_t^{\mathrm{cond}}(\vx|\vz)$；再通过边缘化技巧把它与边缘 score 联系起来。学到 score 之后，可以利用反向 SDE 将简单终点分布重新推回数据分布。训练时虽然真实边缘 score 不可直接计算，但 Denoising Score Matching 定理说明：最小化可计算的条件损失 $\mathcal{L}_{\mathrm{DSM}}$，等价于最小化理想的边缘损失 $\mathcal{L}_{\mathrm{SM}}$。对于 Gaussian path，这个目标进一步退化成对噪声的简单回归问题。
\end{summarybox}
